\documentclass[12pt,a4paper]{exam}
\usepackage[utf8]{inputenc}
\usepackage{amsmath,amssymb,amsfonts}
\usepackage{geometry}
\usepackage{enumitem}
\usepackage{graphicx}
\usepackage{hyperref}
\usepackage{xcolor}
\geometry{a4paper, top=2cm, bottom=2cm, left=2.5cm, right=2.5cm}
\hypersetup{colorlinks=true, linkcolor=blue, urlcolor=blue}
\renewcommand{\thequestion}{\textbf{\arabic{question}}}
\renewcommand{\questionlabel}{\thequestion.}
\pointname{ marks}
\pointsdroppedatright
\bracketedpoints
\setlength{\parindent}{0pt}
\setlength{\emergencystretch}{3em}
\allowdisplaybreaks
\newsavebox{\schmathbox}
\newcommand{\fitmath}[1]{%
  \sbox{\schmathbox}{$\displaystyle #1$}%
  \ifdim\wd\schmathbox>\linewidth
    \resizebox{\linewidth}{!}{\usebox{\schmathbox}}%
  \else
    \usebox{\schmathbox}%
  \fi}

\begin{document}

\thispagestyle{empty}
\begin{center}
  \textbf{\LARGE Diag Solutions} \\[0.3cm]
  \rule{\textwidth}{0.5pt}
\end{center}

\vspace{0.3cm}

\begin{questions}
\question \textbf{1.} Prove that the angle between the two tangents drawn from an external point to a circle is supplementary to the angle subtended by the line segment joining the points of contact at the centre.

\textbf{Answer:}
\begin{center}
\fitmath{\text{Sum is} 180 \text{degrees}}
\end{center}

\textbf{Solution:}
\begin{enumerate}
  \item Let P be the external point and A, B be points of contact.
  \item Consider the quadrilateral OAPB.
  \item Use the property that the sum of angles in a quadrilateral is $360^\circ$ and tangents are perpendicular to radii.
\end{enumerate}
\textbf{Derivation:}
\begin{center}
\fitmath{\begin{aligned}
\angle OAP = 90^\circ, \angle OBP = 90^\circ. \angle APB + \angle OAP + \angle OBP + \angle AOB = 360^\circ \\
\Rightarrow \angle APB + 90^\circ + 90^\circ + \angle AOB = 360^\circ \\
\Rightarrow \angle APB + \angle AOB = 180^\circ.
\end{aligned}}
\end{center}
\hfill \textit{Circles — Tangents to a circle}
\vspace{0.3cm}

\question \textbf{2.} If the circumference of a circle is $176\text{ cm}$, then its radius is:

\textbf{Answer:}
(C)

\textbf{Solution:}
\begin{enumerate}
  \item Use the formula for circumference $C = 2\pi r$.
  \item Substitute $C = 176$ and $\pi = \frac{22}{7}$.
  \item Solve for $r$.
\end{enumerate}
\textbf{Derivation:}
\begin{center}
\fitmath{\begin{aligned}
176 = 2 \times \frac{22}{7} \times r \\
\implies 176 = \frac{44}{7} \times r \\
\implies r = \frac{176 \times 7}{44} = 4 \times 7 = 28\text{ cm}
\end{aligned}}
\end{center}
\hfill \textit{Areas Related to Circles — Circumference of a Circle}
\vspace{0.3cm}

\question \textbf{3.} The area of a circle whose diameter is $14\text{ cm}$ is:

\textbf{Answer:}
(B)

\textbf{Solution:}
\begin{enumerate}
  \item Find the radius $r = \frac{d}{2} = \frac{14}{2} = 7\text{ cm}$.
  \item Use the area formula $A = \pi r^2$.
  \item Substitute $r = 7$ and compute.
\end{enumerate}
\textbf{Derivation:}
\begin{center}
\fitmath{A = \frac{22}{7} \times 7^2 = \frac{22}{7} \times 49 = 22 \times 7 = 154\text{ cm}^2}
\end{center}
\hfill \textit{Areas Related to Circles — Area of a Circle}
\vspace{0.3cm}

\question \textbf{4.} The area of a quadrant of a circle with radius $r$ is:

\textbf{Answer:}
(D)

\textbf{Solution:}
\begin{enumerate}
  \item A quadrant is one-fourth of a circle.
  \item Area of circle = $\pi r^2$.
  \item Area of quadrant = $\frac{1}{4} \pi r^2$.
\end{enumerate}
\textbf{Derivation:}
\begin{center}
\fitmath{\text{Area} = \frac{90^\circ}{360^\circ} \times \pi r^2 = \frac{1}{4} \pi r^2}
\end{center}
\hfill \textit{Areas Related to Circles — Area of a Sector}
\vspace{0.3cm}

\question \textbf{5.} If the perimeter of a semi-circular protractor is $36\text{ cm}$, its diameter is:

\textbf{Answer:}
(A)

\textbf{Solution:}
\begin{enumerate}
  \item Perimeter of a semi-circle is $\pi r + 2r = r(\pi + 2)$.
  \item Set $r(\frac{22}{7} + 2) = 36$.
  \item Solve for $r$ and then diameter $d = 2r$.
\end{enumerate}
\textbf{Derivation:}
\begin{center}
\fitmath{\begin{aligned}
r(\frac{22+14}{7}) = 36 \\
\implies r(\frac{36}{7}) = 36 \\
\implies r = 7. \text{ Thus, } d = 2 \times 7 = 14\text{ cm}.
\end{aligned}}
\end{center}
\hfill \textit{Areas Related to Circles — Perimeter of Semi-circle}
\vspace{0.3cm}

\question \textbf{6.} The area of a sector of a circle with radius $6\text{ cm}$ and central angle $60^\circ$ is:

\textbf{Answer:}
(C)

\textbf{Solution:}
\begin{enumerate}
  \item Use formula $A = \frac{\theta}{360} \times \pi r^2$.
  \item Substitute $\theta = 60$, $r = 6$.
\end{enumerate}
\textbf{Derivation:}
\begin{center}
\fitmath{A = \frac{60}{360} \times \pi \times 6^2 = \frac{1}{6} \times 36\pi = 6\pi\text{ cm}^2}
\end{center}
\hfill \textit{Areas Related to Circles — Area of a Sector}
\vspace{0.3cm}

\question \textbf{7.} Assertion (A): The area of a circle with radius $7$ cm is $154$ cm$^2$. Reason (R): The area of a circle is given by the formula $\pi r^2$.

\textbf{Answer:}
(A)

\textbf{Solution:}
\begin{enumerate}
  \item Calculate area using $\pi r^2$
  \item Area = $(22/7) \times 7 \times 7 = 154$ cm$^2$
  \item Both statements are true and R is the correct explanation
\end{enumerate}
\textbf{Derivation:}
\begin{center}
\fitmath{\text{Area} = \pi r^2 = \frac{22}{7} \times 7^2 = 154}
\end{center}
\hfill \textit{Areas Related to Circles — Area of a Circle}
\vspace{0.3cm}

\question \textbf{8.} Assertion (A): The circumference of a circle is $44$ cm, then its radius is $7$ cm. Reason (R): Circumference of a circle is $2\pi r$.

\textbf{Answer:}
(A)

\textbf{Solution:}
\begin{enumerate}
  \item Use $2\pi r = 44$
  \item $2 \times (22/7) \times r = 44$
  \item $r = (44 \times 7) / 44 = 7$ cm
\end{enumerate}
\textbf{Derivation:}
\begin{center}
\fitmath{\begin{aligned}
2 \times \frac{22}{7} \times r = 44 \\
\implies r = 7
\end{aligned}}
\end{center}
\hfill \textit{Areas Related to Circles — Circumference of a Circle}
\vspace{0.3cm}

\question \textbf{9.} Assertion (A): The area of a semicircle of radius $r$ is $\frac{1}{2}\pi r^2$. Reason (R): A semicircle is half of a circle.

\textbf{Answer:}
(A)

\textbf{Solution:}
\begin{enumerate}
  \item Area of circle is $\pi r^2$
  \item Semicircle is exactly half, so area is $(1/2)\pi r^2$
  \item Both statements are true and R is the correct explanation
\end{enumerate}
\textbf{Derivation:}
Area\_\{semicircle\} = $\frac{1}{2} \times$ Area\_\{circle\} = $\frac{1}{2}\pi r^2$
\hfill \textit{Areas Related to Circles — Area of Semicircle}
\vspace{0.3cm}

\question \textbf{10.} Assertion (A): If the perimeter of a circle is equal to that of a square, then the area of the circle is greater than the area of the square. Reason (R): The area of a circle with a given perimeter is always the maximum possible area for any plane figure.

\textbf{Answer:}
(A)

\textbf{Solution:}
\begin{enumerate}
  \item Let perimeter be $P$. Circle $r = P/2$, Area = $P^2/4$. Square side = $P/4$, Area = $P^2/16$.
  \item Since $4\pi \approx 12.56 < 16$, $P^2/12.56 > P^2/16$.
  \item Both statements are correct.
\end{enumerate}
\textbf{Derivation:}
\begin{center}
\fitmath{A_{\text{circle}} = \frac{P^2}{4\pi} > A_{\text{square}} = \frac{P^2}{16}}
\end{center}
\hfill \textit{Areas Related to Circles — Comparison of Areas}
\vspace{0.3cm}

\question \textbf{11.} Assertion (A): The length of an arc of a sector of a circle with radius $r$ and angle $\theta$ is $\frac{\theta}{360} \times 2\pi r$. Reason (R): The circumference of a circle is $2\pi r$.

\textbf{Answer:}
(A)

\textbf{Solution:}
\begin{enumerate}
  \item An arc is a fraction of the circumference determined by the angle $\theta$.
  \item The fraction is $\theta/360$. Multiply by total circumference $2\pi r$.
  \item Both are true and R explains the basis of the formula.
\end{enumerate}
\textbf{Derivation:}
\begin{center}
\fitmath{\text{Arc Length} = \frac{\theta}{360} \times 2\pi r}
\end{center}
\hfill \textit{Areas Related to Circles — Length of an Arc}
\vspace{0.3cm}

\question \textbf{12.} Find the area of a sector of a circle with radius $6 \text{ cm}$ if the angle of the sector is $60^\circ$. (Use $\pi = \frac{22}{7}$)

\textbf{Answer:}
\begin{center}
\fitmath{\frac{132}{7} \text{ cm}^2}
\end{center}

\textbf{Solution:}
\begin{enumerate}
  \item Identify the formula for the area of a sector as $\frac{\theta}{360^\circ} \times \pi r^2$.
  \item Substitute $r = 6$ and $\theta = 60^\circ$ into the formula and simplify.
\end{enumerate}
\textbf{Derivation:}
\begin{center}
\fitmath{\text{Area} = \frac{60}{360} \times \frac{22}{7} \times 6^2 = \frac{1}{6} \times \frac{22}{7} \times 36 = 6 \times \frac{22}{7} = \frac{132}{7} \text{ cm}^2}
\end{center}
\hfill \textit{Areas Related to Circles — Area of Sector}
\vspace{0.3cm}

\question \textbf{13.} The length of the minute hand of a clock is $14 \text{ cm}$. Find the area swept by the minute hand in $5$ minutes.

\textbf{Answer:}
\begin{center}
\fitmath{\frac{154}{3} \text{ cm}^2}
\end{center}

\textbf{Solution:}
\begin{enumerate}
  \item Calculate the angle swept by the minute hand in 1 minute ($360^\circ/60 = 6^\circ$), so for 5 minutes, $\theta = 30^\circ$.
  \item Apply the area of sector formula with $r = 14$ and $\theta = 30^\circ$.
\end{enumerate}
\textbf{Derivation:}
\begin{center}
\fitmath{\text{Area} = \frac{30}{360} \times \frac{22}{7} \times 14^2 = \frac{1}{12} \times \frac{22}{7} \times 196 = \frac{1}{12} \times 22 \times 28 = \frac{154}{3} \text{ cm}^2}
\end{center}
\hfill \textit{Areas Related to Circles — Application of Sector Area}
\vspace{0.3cm}

\question \textbf{14.} Find the area of a quadrant of a circle whose circumference is $44 \text{ cm}$. (Use $\pi = \frac{22}{7}$)

\textbf{Answer:}
\begin{center}
\fitmath{38.5 \text{ cm}^2}
\end{center}

\textbf{Solution:}
\begin{enumerate}
  \item Use the circumference formula $2\pi r = 44$ to find the radius $r$.
  \item Calculate the area of the quadrant using $\frac{1}{4} \pi r^2$.
\end{enumerate}
\textbf{Derivation:}
\begin{center}
\fitmath{\begin{aligned}
2 \times \frac{22}{7} \times r = 44 \\
\implies r = 7 \text{ cm}. \text{ Area} = \frac{1}{4} \times \frac{22}{7} \times 7^2 = \frac{1}{4} \times 22 \times 7 = 38.5 \text{ cm}^2
\end{aligned}}
\end{center}
\hfill \textit{Areas Related to Circles — Area of Circle Parts}
\vspace{0.3cm}

\question \textbf{15.} If the sum of the areas of two circles with radii $R_1$ and $R_2$ is equal to the area of a circle of radius $R$, prove that $R_1^2 + R_2^2 = R^2$.

\textbf{Answer:}
\begin{center}
\fitmath{R_1^2 + R_2^2 = R^2}
\end{center}

\textbf{Solution:}
\begin{enumerate}
  \item Write the equation based on the given condition: $\pi R_1^2 + \pi R_2^2 = \pi R^2$.
  \item Divide the entire equation by $\text{ }$ to obtain the result.
\end{enumerate}
\textbf{Derivation:}
\begin{center}
\fitmath{\begin{aligned}
\pi R_1^2 + \pi R_2^2 = \pi R^2 \\
\implies \pi(R_1^2 + R_2^2) = \pi R^2 \\
\implies R_1^2 + R_2^2 = R^2
\end{aligned}}
\end{center}
\hfill \textit{Areas Related to Circles — Area of Circle}
\vspace{0.3cm}

\question \textbf{16.} A chord of a circle of radius $10 \text{ cm}$ subtends a right angle at the center. Find the area of the corresponding minor segment. (Use $\pi = 3.14$)

\textbf{Answer:}
\begin{center}
\fitmath{28.5 \text{ cm}^2}
\end{center}

\textbf{Solution:}
\begin{enumerate}
  \item Calculate the area of the sector: $\frac{90}{360} \times \pi r^2$.
  \item Subtract the area of the triangle $\frac{1}{2} r^2 \sin(90^\circ)$ from the sector area.
\end{enumerate}
\textbf{Derivation:}
\begin{center}
\fitmath{\text{Area} = \frac{90}{360} \times 3.14 \times 100 - \frac{1}{2} \times 10 \times 10 = 78.5 - 50 = 28.5 \text{ cm}^2}
\end{center}
\hfill \textit{Areas Related to Circles — Area of Segment}
\vspace{0.3cm}

\question \textbf{17.} Find the area of the shaded region in a circle of radius $14$ cm, where a sector subtends an angle of $60^{\circ}$ at the center. (Use $\pi = \frac{22}{7}$)

\textbf{Answer:}
\begin{center}
\fitmath{102.67 cm^2}
\end{center}

\textbf{Solution:}
\begin{enumerate}
  \item Identify the formula for the area of a sector: $A = \frac{\theta}{360^{\circ}} \times \pi r^2$.
  \item Substitute the given values: $\theta = 60^{\circ}$ and $r = 14$ cm.
  \item Calculate the area: $\frac{60}{360} \times \frac{22}{7} \times 14 \times 14$.
\end{enumerate}
\textbf{Derivation:}
\begin{center}
\fitmath{\text{Area} = \frac{60}{360} \times \frac{22}{7} \times 14 \times 14 = \frac{1}{6} \times 22 \times 2 \times 14 = \frac{616}{6} = 102.67 \text{ cm}^2}
\end{center}
\hfill \textit{Areas Related to Circles — Area of a Sector}
\vspace{0.3cm}

\question \textbf{18.} The length of the minute hand of a clock is $14$ cm. Find the area swept by the minute hand in $5$ minutes.

\textbf{Answer:}
\begin{center}
\fitmath{\frac{154}{3} cm^2}
\end{center}

\textbf{Solution:}
\begin{enumerate}
  \item Determine the angle swept by the minute hand in $60$ minutes ($360^{\circ}$), so in $1$ minute it is $6^{\circ}$.
  \item Calculate the angle for $5$ minutes: $5 \times 6^{\circ} = 30^{\circ}$.
  \item Apply the sector area formula with $r = 14$ cm and $\theta = 30^{\circ}$.
\end{enumerate}
\textbf{Derivation:}
\begin{center}
\fitmath{\text{Area} = \frac{30}{360} \times \frac{22}{7} \times 14 \times 14 = \frac{1}{12} \times 22 \times 2 \times 14 = \frac{616}{12} = \frac{154}{3} \text{ cm}^2}
\end{center}
\hfill \textit{Areas Related to Circles — Applications of Sector Area}
\vspace{0.3cm}

\question \textbf{19.} Find the area of a quadrant of a circle whose circumference is $44$ cm.

\textbf{Answer:}
\begin{center}
\fitmath{38.5 cm^2}
\end{center}

\textbf{Solution:}
\begin{enumerate}
  \item Use the circumference formula $2\pi r = 44$ to find the radius.
  \item Solve for $r$: $r = \frac{44}{2\pi} = \frac{22}{\pi} = 7$ cm.
  \item Calculate the area of the quadrant: $\frac{1}{4} \pi r^2$.
\end{enumerate}
\textbf{Derivation:}
\begin{center}
\fitmath{r = \frac{44 \times 7}{2 \times 22} = 7 \text{ cm}. \text{ Area} = \frac{1}{4} \times \frac{22}{7} \times 7 \times 7 = \frac{154}{4} = 38.5 \text{ cm}^2}
\end{center}
\hfill \textit{Areas Related to Circles — Area of a Quadrant}
\vspace{0.3cm}

\question \textbf{20.} A chord of a circle of radius $10$ cm subtends a right angle at the center. Find the area of the corresponding minor segment. (Use $\pi = 3.14$)

\textbf{Answer:}
\begin{center}
\fitmath{28.5 cm^2}
\end{center}

\textbf{Solution:}
\begin{enumerate}
  \item Find area of the sector: $\frac{90}{360} \times \pi r^2$.
  \item Find area of the triangle formed by the chord and radii: $\frac{1}{2} \times r \times r$.
  \item Subtract the triangle area from the sector area.
\end{enumerate}
\textbf{Derivation:}
Area\_\{sector\} = 0.25 $\times 3.14 \times 100 = 78.5. \text{ Area}_{tri} = 0.5 \times 10 \times 10 = 50. \text{ Area}_{seg} = 78.5 - 50 = 28.5 \text{ cm}^2$
\hfill \textit{Areas Related to Circles — Area of a Segment}
\vspace{0.3cm}

\question \textbf{21.} The radii of two circles are $8$ cm and $6$ cm respectively. Find the radius of the circle having area equal to the sum of the areas of the two circles.

\textbf{Answer:}
\begin{center}
\fitmath{10 cm}
\end{center}

\textbf{Solution:}
\begin{enumerate}
  \item Set up the equation: $\pi R^2 = \pi r_1^2 + \pi r_2^2$.
  \item Simplify to $R^2 = r_1^2 + r_2^2$.
  \item Substitute values $r_1=8, r_2=6$ and solve for $R$.
\end{enumerate}
\textbf{Derivation:}
\begin{center}
\fitmath{R^2 = 8^2 + 6^2 = 64 + 36 = 100. \text{ Therefore, } R = \sqrt{100} = 10 \text{ cm}.}
\end{center}
\hfill \textit{Areas Related to Circles — Areas of Circles}
\vspace{0.3cm}

\question \textbf{22.} A chord $AB$ of a circle of radius $15 \text{ cm}$ subtends an angle of $60^\circ$ at the centre. Find the areas of the corresponding minor and major segments of the circle. (Use $\pi = 3.14$ and $\sqrt{3} = 1.73$)

\textbf{Answer:}
\begin{center}
\fitmath{\text{Area of minor segment} = 20.4375 \text{ cm}^2; \text{Area of major segment} = 686.0625 \text{ cm}^2}
\end{center}

\textbf{Solution:}
\begin{enumerate}
  \item Identify the radius $r = 15$ and central angle $\theta = 60^\circ$.
  \item Calculate the area of sector $OAB = \frac{\theta}{360} \times \pi r^2$.
  \item Calculate the area of $\triangle OAB$. Since $\theta = 60^\circ$ and $OA = OB$, it is an equilateral triangle, so Area $= \frac{\sqrt{3}}{4} r^2$.
  \item Subtract the triangle area from the sector area to find the minor segment.
  \item Subtract the minor segment area from the total circle area to find the major segment.
\end{enumerate}
\textbf{Derivation:}
Area of sector  = $\frac{60}{360} \times 3.14 \times 15^2 = \frac{1}{6} \times 3.14 \times 225 = 117.75 \text{ cm}^2. \text{ Area of } \triangle$ OAB = $\frac{\sqrt{3}}{4} \times 15^2 = 0.4325 \times 225 = 97.3125 \text{ cm}^2. \text{ Minor segment} = 117.75 - 97.3125 = 20.4375 \text{ cm}^2. \text{ Area of circle} = 3.14 \times 15^2 = 706.5 \text{ cm}^2. \text{ Major segment} = 706.5 - 20.4375 = 686.0625 \text{ cm}^2.$ 
\hfill \textit{Areas Related to Circles — Area of Segment}
\vspace{0.3cm}

\question \textbf{23.} An umbrella has 8 ribs which are equally spaced. Assuming the umbrella to be a flat circle of radius $45 \text{ cm}$, find the area between the two consecutive ribs of the umbrella.

\textbf{Answer:}
\begin{center}
\fitmath{\text{Area} = \frac{22275}{28} \text{ cm}^2 \text{or approx} 795.54 \text{ cm}^2}
\end{center}

\textbf{Solution:}
\begin{enumerate}
  \item Determine the total number of sectors formed by 8 ribs, which is 8.
  \item Calculate the central angle of each sector: $\theta = 360^\circ / 8 = 45^\circ$.
  \item Use the formula for the area of a sector: $A = \frac{\theta}{360} \times \pi r^2$.
  \item Substitute $r = 45$ and $\theta = 45^\circ$ into the formula.
  \item Simplify the expression to get the final area.
\end{enumerate}
\textbf{Derivation:}
\begin{center}
\fitmath{\text{Area} = \frac{45}{360} \times \frac{22}{7} \times 45 \times 45 = \frac{1}{8} \times \frac{22}{7} \times 2025 = \frac{11 \times 2025}{4 \times 7} = \frac{22275}{28} \text{ cm}^2.}
\end{center}
\hfill \textit{Areas Related to Circles — Area of Sector}
\vspace{0.3cm}

\question \textbf{24.} Find the area of the shaded region in a square $\text{ABCD}$ of side $14 \text{ cm}$, where four circles are drawn with centers at the four vertices of the square, each having a radius of $7 \text{ cm}$.

\textbf{Answer:}
\begin{center}
\fitmath{\text{Area} = 42 \text{ cm}^2}
\end{center}

\textbf{Solution:}
\begin{enumerate}
  \item Calculate the area of the square $\text{ABCD} = \text{side}^2$.
  \item Identify that each corner forms a quadrant of a circle of radius $7 \text{ cm}$.
  \item Calculate the sum of the areas of the four quadrants, which is equivalent to the area of one full circle of radius $7 \text{ cm}$.
  \item Subtract the area of the circle from the area of the square.
  \item The remaining area is the shaded region.
\end{enumerate}
\textbf{Derivation:}
Area of square  = 14 $\times 14 = 196 \text{ cm}^2. \text{ Area of 4 quadrants} = 4 \times (\frac{1}{4} \pi r^2) = \pi r^2 = \frac{22}{7} \times 7 \times 7 = 154 \text{ cm}^2. \text{ Shaded area} = 196 - 154 = 42 \text{ cm}^2.$ 
\hfill \textit{Areas Related to Circles — Combination of Plane Figures}
\vspace{0.3cm}

\question \textbf{25.} A round table cover has six equal designs as shown in the figure. If the radius of the cover is $28 \text{ cm}$, find the cost of making the designs at the rate of $\text{ \text{Rs.~} } 0.35 \text{ per cm}^2$. (Use $\sqrt{3} = 1.7$)

\textbf{Answer:}
\begin{center}
\fitmath{\text{Cost} = \text{ \text{Rs.~} } 162.68}
\end{center}

\textbf{Solution:}
\begin{enumerate}
  \item Recognize that each of the 6 designs forms a segment of a circle.
  \item Central angle for each segment is $360^\circ / 6 = 60^\circ$.
  \item Calculate the area of one segment: Area of sector - Area of equilateral triangle.
  \item Multiply the area of one segment by 6 to get the total area of the designs.
  \item Multiply the total area by the rate to find the cost.
\end{enumerate}
\textbf{Derivation:}
Area of one segment  = $\frac{60}{360} \times \frac{22}{7} \times 28^2 - \frac{\sqrt{3}}{4} \times 28^2 = \frac{1}{6} \times \frac{22}{7} \times 784 - 1.7 \times 196 = 410.66 - 333.2 = 77.46 \text{ cm}^2. \text{ Total Area} = 6 \times 77.46 = 464.76 \text{ cm}^2. \text{ Cost} = 464.76 \times 0.35 = 162.666 \approx 162.68 \text{ ₹}.$ 
\hfill \textit{Areas Related to Circles — Area of Segment}
\vspace{0.3cm}

\question \textbf{26.} The mean of the first 5 natural numbers is:

\textbf{Answer:}
(B)

\textbf{Solution:}
\begin{enumerate}
  \item Identify the first 5 natural numbers: 1, 2, 3, 4, 5.
  \item Sum the numbers: 1 + 2 + 3 + 4 + 5 = 15.
  \item Divide the sum by the count of numbers: 15 / 5 = 3.
\end{enumerate}
\textbf{Derivation:}
\begin{center}
\fitmath{\bar{x} = \frac{1 + 2 + 3 + 4 + 5}{5} = \frac{15}{5} = 3}
\end{center}
\hfill \textit{Statistics — Mean}
\vspace{0.3cm}

\question \textbf{27.} Find the mode of the following data: 2, 3, 4, 3, 5, 3, 6, 3.

\textbf{Answer:}
(C)

\textbf{Solution:}
\begin{enumerate}
  \item Observe the frequency of each number.
  \item The number 3 appears four times, which is more frequent than any other number.
  \item Therefore, the mode is 3.
\end{enumerate}
\textbf{Derivation:}
\begin{center}
\fitmath{\text{Mode is the observation with the highest frequency. Since 3 occurs 4 times, Mode = 3.}}
\end{center}
\hfill \textit{Statistics — Mode}
\vspace{0.3cm}

\question \textbf{28.} If the mean of $x, x+2, x+4, x+6, x+8$ is 24, find the value of $x$.

\textbf{Answer:}
(A)

\textbf{Solution:}
\begin{enumerate}
  \item Sum the observations: x + (x+2) + (x+4) + (x+6) + (x+8) = 5x + 20.
  \item Divide by number of observations (5) and equate to mean (24).
  \item Solve for x: (5x + 20) / 5 = 24 => x + 4 = 24 => x = 20.
\end{enumerate}
\textbf{Derivation:}
\begin{center}
\fitmath{\begin{aligned}
\frac{5x + 20}{5} = 24 \\
\implies x + 4 = 24 \\
\implies x = 20
\end{aligned}}
\end{center}
\hfill \textit{Statistics — Mean}
\vspace{0.3cm}

\question \textbf{29.} The median of first 7 prime numbers is:

\textbf{Answer:}
(B)

\textbf{Solution:}
\begin{enumerate}
  \item List the first 7 prime numbers: 2, 3, 5, 7, 11, 13, 17.
  \item Since n=7 (odd), the median is the ((n+1)/2)th observation.
  \item The 4th observation is 7.
\end{enumerate}
\textbf{Derivation:}
\begin{center}
\fitmath{\text{Median} = \left(\frac{7+1}{2}\right)\text{th term} = \text{4th term} = 7}
\end{center}
\hfill \textit{Statistics — Median}
\vspace{0.3cm}

\question \textbf{30.} For a given distribution, if Mean = 20 and Mode = 20, find the Median using the empirical relationship.

\textbf{Answer:}
(C)

\textbf{Solution:}
\begin{enumerate}
  \item Use the empirical formula: Mode = 3(Median) - 2(Mean).
  \item Substitute the given values: 20 = 3(Median) - 2(20).
  \item \(
20 = 3(\text{Median}) - 40 => 60 = 3(\text{Median}) => \text{Median} = 20.
\)
\end{enumerate}
\textbf{Derivation:}
\begin{center}
\fitmath{\begin{aligned}
20 = 3(\text{Median}) - 2(20) \\
\implies 20 + 40 = 3(\text{Median}) \\
\implies \text{Median} = 60/3 = 20
\end{aligned}}
\end{center}
\hfill \textit{Statistics — Empirical Relationship}
\vspace{0.3cm}

\question \textbf{31.} Assertion (A): The mean of the first five natural numbers is 3. Reason (R): The mean of a set of observations is the sum of observations divided by the total number of observations.

\textbf{Answer:}
(A)

\textbf{Solution:}
\begin{enumerate}
  \item \(
\text{Sum of first} 5 \text{natural numbers} = 1 + 2 + 3 + 4 + 5 = 15
\)
  \item \(
\text{Mean} = Sum / \text{Number of observations} = 15 / 5 = 3
\)
  \item Assertion is true and Reason is the correct definition/formula.
\end{enumerate}
\textbf{Derivation:}
\begin{center}
\fitmath{\text{Mean} = \frac{\sum x_i}{n} = \frac{1+2+3+4+5}{5} = \frac{15}{5} = 3}
\end{center}
\hfill \textit{Statistics — Mean of ungrouped data}
\vspace{0.3cm}

\question \textbf{32.} Assertion (A): The empirical relationship between mean, median, and mode is $3 \text{Median} = \text{Mode} + 2 \text{Mean}$. Reason (R): For a symmetric distribution, mean, median, and mode are equal.

\textbf{Answer:}
(A)

\textbf{Solution:}
\begin{enumerate}
  \item The empirical formula is indeed 3 Median = Mode + 2 Mean.
  \item In a perfectly symmetric distribution, mean = median = mode, which satisfies the formula.
  \item Both statements are true and related to the central tendency properties.
\end{enumerate}
\textbf{Derivation:}
\begin{center}
\fitmath{3 \text{Median} = \text{Mode} + 2 \text{Mean}}
\end{center}
\hfill \textit{Statistics — Relationship between mean, median and mode}
\vspace{0.3cm}

\question \textbf{33.} Assertion (A): The mode of the data 2, 3, 4, 3, 2, 3 is 3. Reason (R): Mode is the observation that occurs most frequently.

\textbf{Answer:}
(A)

\textbf{Solution:}
\begin{enumerate}
  \item Count frequencies: 2 appears twice, 3 appears three times, 4 appears once.
  \item The most frequent value is 3.
  \item Assertion is true and Reason is the correct definition of mode.
\end{enumerate}
\textbf{Derivation:}
\begin{center}
\fitmath{\text{Frequency of 2} = 2, \text{Frequency of 3} = 3, \text{Frequency of 4} = 1. \therefore \text{Mode} = 3}
\end{center}
\hfill \textit{Statistics — Mode of ungrouped data}
\vspace{0.3cm}

\question \textbf{34.} Assertion (A): If the mean of a distribution is 20 and the sum of observations is 100, then the number of observations is 5. Reason (R): The sum of observations = Mean $\times$ Number of observations.

\textbf{Answer:}
(A)

\textbf{Solution:}
\begin{enumerate}
  \item \(
\text{Given Mean} = 20, Sum = 100.
\)
  \item Using formula: 100 = 20 * n => n = 100 / 20 = 5.
  \item Assertion is true and Reason provides the correct formula used for calculation.
\end{enumerate}
\textbf{Derivation:}
\begin{center}
\fitmath{n = \frac{\sum x_i}{\bar{x}} = \frac{100}{20} = 5}
\end{center}
\hfill \textit{Statistics — Mean of ungrouped data}
\vspace{0.3cm}

\question \textbf{35.} Assertion (A): The median of the data 5, 8, 12, 16, 20 is 12. Reason (R): The median is always the middle value when the data is arranged in ascending order.

\textbf{Answer:}
(A)

\textbf{Solution:}
\begin{enumerate}
  \item Data: 5, 8, 12, 16, 20 (already in ascending order).
  \item There are 5 observations (odd). Median is the (n+1)/2 th term.
  \item 3rd term is 12. Assertion is true and Reason is the correct definition.
\end{enumerate}
\textbf{Derivation:}
\begin{center}
\fitmath{\text{Median} = \left(\frac{n+1}{2}\right)^{\text{th}} \text{term} = \left(\frac{5+1}{2}\right)^{\text{th}} = 3^{\text{rd}} \text{term} = 12}
\end{center}
\hfill \textit{Statistics — Median of ungrouped data}
\vspace{0.3cm}

\question \textbf{36.} The mean of 5 observations is 15. If each observation is increased by 3, find the new mean.

\textbf{Answer:}
\begin{center}
\fitmath{18}
\end{center}

\textbf{Solution:}
\begin{enumerate}
  \item Use the property that if each observation is increased by a constant k, the mean also increases by k.
  \item Add the constant to the original mean.
\end{enumerate}
\textbf{Derivation:}
\begin{center}
\fitmath{\text{Original Mean } (\bar{x}) = 15. \text{ New Mean } = \bar{x} + 3 = 15 + 3 = 18.}
\end{center}
\hfill \textit{Statistics — Mean of grouped/ungrouped data}
\vspace{0.3cm}

\question \textbf{37.} Find the class mark of the class interval 150-200.

\textbf{Answer:}
\begin{center}
\fitmath{175}
\end{center}

\textbf{Solution:}
\begin{enumerate}
  \item Identify the lower limit (L) and upper limit (U) of the interval.
  \item Apply the formula: Class mark = (L + U) / 2.
\end{enumerate}
\textbf{Derivation:}
\begin{center}
\fitmath{\text{Class mark} = \frac{150 + 200}{2} = \frac{350}{2} = 175.}
\end{center}
\hfill \textit{Statistics — Class intervals and class marks}
\vspace{0.3cm}

\question \textbf{38.} If the mode of a data set is 45 and the mean is 27, find the median using the empirical relationship.

\textbf{Answer:}
\begin{center}
\fitmath{33}
\end{center}

\textbf{Solution:}
\begin{enumerate}
  \item State the empirical formula: Mode = 3 Median - 2 Mean.
  \item Substitute the given values and solve for Median.
\end{enumerate}
\textbf{Derivation:}
\begin{center}
\fitmath{\begin{aligned}
45 = 3(\text{Median}) - 2(27) \\
\implies 45 = 3(\text{Median}) - 54 \\
\implies 99 = 3(\text{Median}) \\
\implies \text{Median} = 33.
\end{aligned}}
\end{center}
\hfill \textit{Statistics — Relationship between Mean, Median and Mode}
\vspace{0.3cm}

\question \textbf{39.} The frequency of a class is 12 and the total number of observations is 40. Find the relative frequency of this class.

\textbf{Answer:}
\begin{center}
\fitmath{0.3}
\end{center}

\textbf{Solution:}
\begin{enumerate}
  \item Use the formula for relative frequency: (Frequency of class) / (Total number of observations).
  \item Calculate the division.
\end{enumerate}
\textbf{Derivation:}
\begin{center}
\fitmath{\text{Relative Frequency} = \frac{12}{40} = \frac{3}{10} = 0.3.}
\end{center}
\hfill \textit{Statistics — Frequency distribution}
\vspace{0.3cm}

\question \textbf{40.} Given the data: 12, 15, 18, 12, 15, 12, 18. Find the range of the data.

\textbf{Answer:}
\begin{center}
\fitmath{6}
\end{center}

\textbf{Solution:}
\begin{enumerate}
  \item Identify the maximum and minimum values in the data set.
  \item Subtract the minimum from the maximum.
\end{enumerate}
\textbf{Derivation:}
\begin{center}
\fitmath{\text{Max} = 18, \text{Min} = 12. \text{Range} = 18 - 12 = 6.}
\end{center}
\hfill \textit{Statistics — Measures of dispersion}
\vspace{0.3cm}

\end{questions}

\end{document}
